% !TEX root = ../main.tex
\pagebreak
\lecture{3}{2026-11-12}{Anekdote der politischen/modernen Oekonomie}

\section{Anekdote der modernen ökonomischen Wissenschaft}

Stellen wir uns einen Zauberer vor. Ihn nennen wir Robinson.
Robinson hatte einen unterhaltsamen neuen Trick für seine
nächste Zaubershow entwickelt: eine Maschine, die Äpfel in
Goldstücke verwandelte. Die Stadt war zunächst voller
Werbeplakate, und an den Mittagstischen wurde nur noch über
seinen nächsten Auftritt gesprochen, wie es bei einer solchen
Versprechung zu erwarten war. Beim erstmaligen Vorführen dieses
Tricks konnte man die Spannung im Saal spüren; alle waren
still, blickten nach vorn und waren erstaunt, als 5 Kilo Äpfel,
die vom Förderband in die Maschine gebracht wurden, sich
plötzlich auf der anderen Seite in reines Gold verwandelten.
Das Bild des glänzenden gelben Haufens spiegelte sich in den
Pupillen aller Zuschauer und in ihren ruhelosen Gedanken, als
sie am Ende des Showabends entsetzt nach Hause fuhren. In den
darauffolgenden Tagen entwickelten Laien der Zauberkunst
allerlei Theorien über die Funktionsweise der – nach den einen
– revolutionären neuen Technologie oder geschickten Illusion –
nach den anderen. Als auch nach mehreren Wochen die aufgeheizte
Diskussion unter den Bürgern des Landes noch keine Ruhe
gefunden hatte, erreichte die Nachricht auch die entlegene Welt
der Akademiker, und Herr Prof. h.c. Dr. Marsbart erklärte der
ungeschulten Öffentlichkeit, er sei dabei, das grosse Mysterium
aufzulösen.

Er begann also seine Arbeit im weiten Ozean seiner Bibliothek.
Dort wählte er eine passende Methodologie aus, suchte nach
ähnlichen Erfahrungen in anderen empirischen Studien und fühlte
sich nach einigen Tagen überzeugt von seinem eigenen
wissenschaftlichen Untersuchungsvermögen, das alchemische
Geheimnis des Ganzen zu enthüllen. Der Professor kontaktierte
Robinson, um ihn um Erlaubnis zu bitten, seine Maschine zu
untersuchen. Robinson akzeptierte und zeigte dabei
erstaunlicherweise keinerlei Zeichen von Unsicherheit oder
Angst um seine seit Kurzem steil aufsteigende Karriere. Er
hatte übrigens auch seine Recherche über Herrn Dr. h.c.
Marsbart getrieben.

Nach drei Reisetagen war Marsbart angekommen, und ohne sich
viel Zeit für soziale Kleinigkeiten zu nehmen, setzte er sich
sofort an die Maschine, und mit Stift und Papier begann er,
seine Beobachtungen zu protokollieren.

Mehrere Monate dauerte sein Vorhaben, während die
Öffentlichkeit gierig auf Antworten wartete: Kinder fragten
ihre Eltern aus, Eltern verloren deswegen ihren Schlaf,
Esoteriker und Betrunkene debattierten nächtelang in der Bar;
aber die Zeit war endlich gekommen, an dem Herr Dr. h.c.
Marsbart seine wissenschaftliche Auflösung der grossen Frage
veröffentlichen würde.

Neugierige Augen besetzten Stühle und Treppen des grössten
Universitätssaals. Die vordersten Plätze waren für die Presse
reserviert, um eine möglichst genaue Abbildung der kommenden
Fakten zu dokumentieren. Auch am Universitätsplatz sammelten
sich Menschen aller Demografien, um die rudimentäre Übertragung
der Worte des grossen Wissenschaftlers zu verfolgen. Hinten,
ganz unauffällig verkleidet, befand sich der Zauberer und
wartete geduldig mit einem schlauen Lächeln auf die Rede. Der
Moment war also endlich gekommen. Prof. h.c. Dr. Marsbart
stellte sich hinter das Podium, in seinem besten Anzug und mit
seiner Lieblingskrawatte, und begann, seine neu gewonnenen
Kenntnisse der Fakten zu vermitteln.

„Meine Damen und Herren, Liebhaber der Wissenschaft und der
Wahrheit, ich trete heute mit grossem Stolz vor Sie hin, da mir
die Wahrheit der Sache klar geworden ist.“ Alle klatschten so
lange und enthusiastisch, dass Marsbart sie unterbrechen
musste. „Sie wissen ja alle, um welche Frage es sich hier
handelt, also bedarf sie keiner weiteren Erläuterung.“ Die
Zuschauer waren jetzt still und verfolgten jedes weitere Wort
mit ihrer gesamten Aufmerksamkeit.

„Ich habe jeden Millimeter dieser Maschine unter den
verschiedensten Bedingungen gründlich analysiert, und obwohl
ich Ihnen die ausführlichen Erklärungen meiner exakten
Methodologie erspare, werde ich Ihnen nur Folgendes dazu sagen:
Ihr alchemisches, metaphysisches Verhalten habe ich spitzfindig
systematisiert, hermeneutisch interpretiert, ontologisch
entschlüsselt und vor allem empirisch quantitativ beobachtet.
Der Rest meiner Rede wird sich den Resultaten dieser
methodischen Recherche widmen.“

Das Publikum bewältigte seine momentane Verwirrung, und etwas
genervt hörte es wieder absolut aufmerksam zu.

„Meine erste Erkenntnis bestand darin, dass die Funktionalität
der Maschine unabhängig von den Wetterbedingungen war, sodass
sie bei Sonne, Regen und Schnee in ihrer äpfelverwandelnden
Tätigkeit ungestört weiterarbeitete. Hier kommt jedoch das
interessante Faktum: Diese Wettertoleranz basiert auf folgender
Tatsache: Die Mächtigkeit ihrer Erzeugnisse hängt nur vom
UV-Lichteinfall ab, der von troposphärischen Ereignissen auch
unabhängig ist! Ersteres nenne ich die 'Invariabilität' und
Letzteres die 'photoaureatische Proportionalität'.“

Von den Wortschatzkenntnissen des Professors beeindruckt,
jedoch von der Relevanz seiner Beobachtungen nicht besonders
überzeugt, hörten sie weiter zu.

„Darüber hinaus zeigten meine statistischen Messungen mit
promissorischen p-Werten, dass je reifer die zuzuführende
Frucht, desto höher die Reinheit des goldigen Produkts war,
jeweils in molarer Konzentration von Ammoniak und
Schwefelwasserstoff und in Karat gemessen. Jedoch besteht
dieser Zusammenhang nicht einfach aus einem linearen oder
arithmetischen Term, weil auch der verfaulteste Apfel immerhin
Gold lieferte, auch wenn nur in sehr geringer Konzentration.
Die reifste Frucht erzeugt umgekehrt auch keine 100\% Reinheit.
Dies nenne ich das 'sapio'sche Kovexitätsgesetz'.

In Verbindung mit dem uns allen bekannten Energie- und
Massenerhaltungsgesetz konnte ich feststellen, dass die gesamte
Materie der Zufuhr stets entweder in Gold verwandelt wurde oder
als Unreinheit im metallischen Produkt weiterbestand. Die
Maschine respektiert also die thermodynamischen Gesetze der
Konstanz.

Diese Zusammenhänge erlauben es uns, eine mehrdimensionale
Funktion aufzustellen, die die Masse und Reife der Edukte sowie
den UV-Lichteinfall mit Goldreinheit und Produktmasse
verknüpft.  Diese Funktion habe ich – wenn Sie es mir erlauben
– die 'Marsbart-Funktion' benannt. Ihre mathematischen
Eigenschaften sind hier nicht weiter zu beleuchten, da uns hier
die Suche nach der optimalen Gleichgewichtslage ihrer
Produktivität interessiert. Dazu ist neben der neu entdeckten
Marsbart-Funktion eine Bedingung notwendig. Diese ergibt sich
aus der Tatsache, dass die Maschine kein Perpetuum mobile ist,
sondern ihre Leistung aus einer externen Energiequelle bezieht,
die wiederum nicht unerschöpflich und in ihrer Mächtigkeit
zeitlich limitiert ist.

Vor uns befinden sich jetzt alle notwendigen Elemente für die
Auflösung des Problems: Wenn wir unter der energetischen
Bedingung die Funktion maximieren, kommen wir so auf die
optimale Menge an Äpfeln, auf ihren entsprechenden Reifegrad
und die UV-Lichtintensität, die uns mit den grössten
Gold-Erzeugnissen beliefert. Dies, meine Damen und Herren,
Liebhaber der Wissenschaft, ist mein endgültiger Beitrag zu
dieser ungeheuren Polemik, der das Erklärungsvermögen der
Wissenschaft ein für alle Mal beweist. Vielen Dank für Ihre
Aufmerksamkeit!“

Platz und Saal wären still gewesen, gäbe es nicht den langen
stehenden Applaus der anderen anwesenden Gelehrten, Apfelbauern
und Schmuckhersteller sowie das Seufzen der
Goldmarktspekulanten. Diese waren jedoch genug, um das laute
Lachen des Zauberers und die unzufriedenen, verwirrten
Gesichter der Masse zu dämpfen. Nach dem Vorfall wurde Marsbart
zum Professor ernannt und zum zweifachen Dr. h.c. promoviert;
zudem erhielt er einen Lehrstuhl an der angesehensten
Universität des Landes sowie eine gemütliche technische
Führungsposition in der durch seine Beiträge neu erfundenen
Industrie der bio-metallischen Verwandlung. Der Zauberer
Robinson, zufrieden mit dem Ausgang, den das Ganze genommen
hatte, verfolgte weiterhin seine Leidenschaft für weitere
unterhaltsame Erfindungen und genoss die Show-Einkünfte, die er
sich während dieser ganzen Geschichte erworben hatte. Der
breiten Mehrheit des Volkes wurde das Thema zur blossen
Trivialität, wie es bei jeder revolutionären Erfindung der Fall
ist. In den Bars hatten Esoteriker und Betrunkene weiterhin
argumentativ gestritten, wobei die Betrunkenen schliesslich
durch Aufgabe gewannen.
\pagebreak
\subsection{O professor e o mágico}
Imaginemos um mágico. Vamos chamá-lo de Robinson. Robinson
tinha desenvolvido um novo pequeno truque para o seu
próximo show de mágica: uma máquina que transformava maçãs em
ouro. A cidade primeiro encheu-se de cartazes
publicitários, e na hora do almoço só se falava da sua próxima
apresentação, como era de se esperar de uma promessa dessas. Na
primeira demonstração do truque, dava para sentir a tensão na
plateia; todos estavam calados, olhando para a frente,
espantados quando 5 quilos de maçãs, trazidas pela esteira para
dentro da máquina, apareceram de repente do outro lado
transformadas em ouro puro. A imagem do monte amarelo brilhante
refletia-se nas pupilas de todos os espectadores e nos seus
pensamentos inquietos, quando, no fim da noite de show,
voltaram para casa apavorados. Nos dias seguintes, os leigos da
arte mágica desenvolveram todo tipo de teorias sobre o
funcionamento daquela – segundo uns, revolucionária nova
tecnologia; segundo outros, hábil ilusão. Quando, passadas
várias semanas, a discussão acalorada ainda não tinha
encontrado sossego entre os cidadãos do país, a notícia chegou
também ao remoto mundo dos acadêmicos, e o Sr. Prof. h.c. Dr.
Jegues declarou ao público inculto que estava prestes a
desvendar o grande mistério.

Começou então o seu trabalho no vasto oceano da sua biblioteca.
Lá escolheu uma metodologia adequada, procurou experiências
semelhantes em outros estudos empíricos e, ao fim de alguns
dias, sentiu-se convencido do seu próprio poder de investigação
científica para revelar o segredo alquímico de tudo aquilo. O
professor contatou Robinson para lhe pedir autorização para
examinar a sua máquina. Robinson aceitou e, surpreendentemente,
não mostrou nenhum sinal de insegurança ou medo pela sua
carreira, que ultimamente subia a pique. Aliás, ele também
tinha feito a sua pesquisa sobre o senhor doutor.

Ao fim de três dias de viagem, Jegues chegou e, sem perder
muito tempo com pequenas conveniências sociais, sentou-se logo
junto à máquina e, com lápis e papel, começou a registrar as
suas observações.

O seu projeto durou vários meses, enquanto o público esperava
avidamente por respostas: as crianças interrogavam os pais, os
pais perdiam o sono por causa disso, esotéricos e bêbados
discutiam noites a fio no bar; mas tinha finalmente chegado o
momento em que o Sr. Dr. h.c. publicaria a sua resolução
científica da grande questão.

Olhos curiosos ocupavam cadeiras e escadas do maior auditório
da universidade. Os lugares da frente estavam reservados à
imprensa, para documentar o mais fielmente possível a imagem
dos fatos vindouros. Também na praça da universidade se reuniam
pessoas de todas as demografias para acompanhar a transmissão
rudimentar das palavras do grande cientista. Lá atrás,
totalmente disfarçado e discreto, estava o mágico, à espera
paciente, com um sorriso astuto, do discurso. O momento tinha,
portanto, finalmente chegado. O Prof. h.c. Dr. Jegues
colocou-se atrás do pódio, no seu melhor terno e com a sua
gravata preferida, e começou a transmitir os conhecimentos dos
fatos que acabara de adquirir.

„Minhas senhoras e meus senhores, amantes da ciência e da
verdade, hoje me apresento diante de vocês com grande orgulho,
pois a verdade da coisa tornou-se clara para mim.“ Todos
aplaudiram durante tanto tempo e com tanto entusiasmo que
Jegues teve de interrompê-los. „Todos sabem, naturalmente, de
que questão se trata aqui, portanto ela não necessita de mais
explicações.“ Os espectadores ficaram agora calados e seguiram
cada palavra seguinte com toda a sua atenção.

„Analisei minuciosamente cada milímetro desta máquina sob as
mais diversas condições e, embora poupe vocês das explicações
detalhadas da minha metodologia exata, dir-vos-ei apenas o
seguinte. Isto deverá bastar para convencer qualquer leigo da
cientificidade da minha investigação: sistematizei argutamente
o seu comportamento alquímico e metafísico, interpretei-o
hermeneuticamente, decifrei-o ontologicamente e, sobretudo,
observei-o empírica e quantitativamente. O resto do meu
discurso será dedicado aos resultados desta investigação
metódica.“

O público superou a sua momentânea confusão e, um pouco
irritado, voltou a ouvir com absoluta atenção.

„A minha primeira constatação foi que a funcionalidade da
máquina era independente das condições meteorológicas, de modo
que, com sol, chuva ou neve, continuava sem perturbação a sua
atividade transformadora de maçãs. No entanto, aqui está o fato
interessante: esta tolerância ao tempo baseia-se na seguinte
circunstância: a pureza dos seus produtos depende da
incidência de luz ultravioleta! Ao primeiro chamo a
‚invariabilidade‘ e ao segundo a ‚proporcionalidade
fotoáurea‘.“

Impressionados com o conhecimento vocabular do professor, mas
não particularmente convencidos da relevância das suas
observações, continuaram a ouvir.

„Além disso, as minhas medições estatísticas com valores-p
promissórios mostraram que, quanto mais madura a fruta a ser
fornecida, tanto maior a pureza do produto dourado, medidas
respectivamente em concentração molar de amoníaco e sulfeto de
hidrogênio, e em quilates. No entanto, esta relação não
consiste simplesmente num termo linear ou aritmético, porque
mesmo a maçã mais podre fornecia sempre algum ouro, ainda que em
concentração muito reduzida. Inversamente, a fruta mais madura
também não produz 100\% de pureza. A isto chamo a ‚lei da
sapio-covexidade‘.

Além disso, em ligação com a lei da conservação da energia e da
massa, conhecida por todos nós, pude constatar que toda a
matéria do fornecimento era sempre ou transformada em ouro ou
continuava a existir como impureza no produto metálico. A
máquina respeita, portanto, as leis termodinâmicas da
constância.

Estas relações permitem-nos estabelecer uma função
multidimensional que liga a massa e a maturidade dos edutos,
bem como a incidência de luz UV, à pureza do ouro e à massa do
produto. A esta função chamei – se me permitem – a ‚função de
Jegues‘. As suas propriedades matemáticas não precisam de ser
aqui mais iluminadas, pois o que nos interessa aqui é a procura
da posição de equilíbrio ótima da sua produtividade. Para isso,
além da recém-descoberta função de Jegues, é necessária uma
condição. Esta decorre do fato de a máquina retirar o seu
rendimento de uma fonte de energia externa, que, por sua vez,
não é inesgotável e está limitada no tempo quanto à sua
magnitude.

Diante de nós encontram-se agora todos os elementos necessários
para a resolução do problema: se, sob a condição energética,
maximizarmos a função, chegamos assim à quantidade ótima de
maçãs, ao seu correspondente grau de maturidade e à intensidade
de luz UV que nos fornece os maiores produtos de ouro. Isto,
minhas senhoras e meus senhores, amantes da ciência, é o meu
contributo definitivo para esta enorme polêmica, que prova de
uma vez por todas a capacidade explicativa da ciência. Muito
obrigado pela vossa atenção!“

A praça e auditório estariam silenciosos, não fosse o longo
aplauso de pé dos outros presentes – eruditos, produtores de
maçãs, fabricantes de joias – e os suspiros dos especuladores
do mercado do ouro. Estes bastaram, porém, para abafar a
gargalhada do mágico e os rostos insatisfeitos e confusos da
multidão. Após o causo, Jegues foi promovido a professor
e a duplo Dr. h.c., e recebeu uma cadeira na universidade mais
prestigiada do país, além de uma confortável posição de direcão
técnica na indústria da transformação bio-metálica,
recém-inventada graças aos suas contribuicões. O mágico Robinson,
satisfeito com o desfecho que tudo tinha tomado, continuou a
sua paixão por outras invenções divertidas e aproveitou os
rendimentos do show que tinha adquirido ao longo de toda esta
história. 

Para a maioria, o tema tornou-se numa mera
trivialidade, como acontece com qualquer invenção
revolucionária. Nos bares, esotéricos e bêbados tinham continuado a
debater e por fim os bêbados acabaram por ganhar por desistência.“